\section{Experiments}

\vspace{-1.2ex}
We evaluate whether \methodname improves personalized presentation generation
and multi-turn localized revision. The main text reports personalization results
on a controlled multi-persona, multi-intent user profile bank, a
DeepPresenter-style general-quality check on the same decks, and a diagnostic
matched-pair modify evaluation for localized revision. Protocol details,
baseline conditions, and profile-bank construction are provided in
Appendices~\ref{app:evaluation_protocol_details},
\ref{app:baseline_provenance}, \ref{app:profile_bank_construction},
and~\ref{app:deeppresenter_quality_eval}, with protocol, profile-bank, and
compute summaries in Appendix Tables~\ref{tab:appendix_eval_protocol_summary},
\ref{tab:appendix_profile_library_catalog}, and
\ref{tab:appendix_compute_accounting}; table captions define the reported
metrics.

\subsection{Main Results}

\subsubsection{Personalization}

\begin{table}[t]
\centering
\small
\setlength{\tabcolsep}{6pt}
\renewcommand{\arraystretch}{1.0}
\caption{Persona-alignment judgments for first-pass generation. Scores are
averaged over three personas on a 0--10 scale; higher is better. Bold marks the
best score per metric.}
\label{tab:profile_memory_v6_bestof_main}
\begin{tabular}{clcccc}
\toprule
Framework & Model & Content $\uparrow$ & Structure $\uparrow$ & Visual $\uparrow$ & Specificity $\uparrow$ \\
\midrule
\multirow{3}{*}{\centering DeepPresenter} & GPT-5 & 6.22 & 7.56 & 5.76 & 5.89 \\
& GLM-5 & 6.67 & 7.61 & 5.28 & 7.22 \\
& Gemini 3.1 Pro & 6.89 & 8.00 & 6.78 & 7.44 \\
\midrule
\multirow{3}{*}{\centering SlideTailor} & GPT-5 & 6.78 & 6.00 & 6.39 & 6.33 \\
& GLM-5 & 4.44 & 4.89 & 4.00 & 3.89 \\
& Gemini 3.1 Pro & 4.48 & 5.00 & 4.03 & 4.67 \\
\midrule
\multirow{3}{*}{\centering MemSlides (Ours)} & GPT-5 & 7.11 & 7.33 & 6.00 & 6.67 \\
& GLM-5 & \textbf{9.00} & \textbf{8.78} & \textbf{8.56} & \textbf{8.89} \\
& Gemini 3.1 Pro & 7.77 & 8.64 & 8.24 & 8.56 \\
\bottomrule
\end{tabular}
\end{table}


User profile memory improves round-0 persona alignment across multiple
dimensions. Under the persona-alignment judgments, \textsc{Ours} achieves
all-column wins over both baselines on GLM-5 and Gemini 3.1 Pro. With GPT-5, it
remains ahead of SlideTailor on Content, Structure, and Specificity, and ahead
of DeepPresenter on Content, Visual, and Specificity, while DeepPresenter has
slightly higher Structure (Table~\ref{tab:profile_memory_v6_bestof_main};
Appendix Table~\ref{tab:profile_memory_v6_bestof_gpt_appendix} gives
per-persona GPT-5 details).

\begin{table}[t]
\centering
\small
\setlength{\tabcolsep}{5.2pt}
\renewcommand{\arraystretch}{1.0}
\caption{General-quality evaluation on the
shared three-profile suite.  Constraint, Content, Style, and Avg. use a 1--5 scale;
Diversity is a suite-level normalized DINOv2-Vendi score. The rows compare
independent generated decks from DeepPresenter, SlideTailor, and \methodname
under the same evaluation protocol.}
\label{tab:profile_memory_v6_ppteval_main}
\begin{tabular*}{\textwidth}{@{\extracolsep{\fill}}clccccc@{}}
\toprule
Framework & Model & Constraint $\uparrow$ & Content $\uparrow$ & Style $\uparrow$ & Avg. $\uparrow$ & Diversity $\uparrow$ \\
\midrule
\multirow{3}{*}{\centering DeepPresenter}
& GPT-5 & 4.83 & 3.50 & 3.63 & 3.99 & 0.387 \\
& Gemini 3.1 Pro & 4.17 & 3.33 & 4.00 & 3.83 & 0.370 \\
& GLM-5 & 4.00 & 3.57 & 4.00 & 3.86 & 0.366 \\
\midrule
\multirow{3}{*}{\centering SlideTailor}
& GPT-5 & 3.83 & 2.93 & 4.03 & 3.60 & 0.399 \\
& Gemini 3.1 Pro & 3.83 & 3.20 & 4.00 & 3.68 & 0.364 \\
& GLM-5 & 3.83 & 2.97 & 4.00 & 3.60 & 0.348 \\
\midrule
\multirow{3}{*}{\centering MemSlides ({Ours})}
& GPT-5 & \textbf{5.00} & \textbf{3.60} & 3.90 & \textbf{4.17} & 0.380 \\
& Gemini 3.1 Pro & 3.33 & 3.37 & \textbf{4.10} & 3.60 & \textbf{0.463} \\
& GLM-5 & 3.83 & 3.34 & 4.03 & 3.74 & 0.391 \\
\bottomrule
\end{tabular*}
\end{table}


\begin{table}[t]
\centering
\small
\setlength{\tabcolsep}{6pt}
\renewcommand{\arraystretch}{0.9}
\newcommand{\toolmemorymultiheader}[1]{\shortstack[c]{#1}}
\providecommand{\cmark}{\textcolor{green!50!black}{\checkmark}}
\providecommand{\xmark}{\textcolor{red!70!black}{\(\times\)}}
\newcommand{\toolmemoryyes}{\cmark}
\newcommand{\toolmemoryno}{\xmark}
\caption{Tool-memory ablation on nine diagnostic modify pairs. Completion and
verification are higher-is-better; time and ratio are lower-is-better.}
\label{tab:tool_memory_table1_redesign}
\begin{tabular*}{\textwidth}{@{\extracolsep{\fill}}lccccc@{}}
\toprule
Model & \toolmemorymultiheader{Memory\\Injected} &
\toolmemorymultiheader{Closed-Loop\\Completion $\uparrow$} &
\toolmemorymultiheader{Strict\\Verify $\uparrow$} &
\toolmemorymultiheader{First Correct\\Edit (s) $\downarrow$} &
\toolmemorymultiheader{Core Tool\\Time Ratio $\downarrow$} \\
\midrule
\multirow{2}{*}{GPT-5} & \toolmemoryyes{} & \textbf{1.000} & \textbf{0.646} & \textbf{211.3} & \textbf{0.740$\times$} \\
 & \toolmemoryno{} & 0.667 & 0.294 & 234.2 & 1.000$\times$ \\
\midrule
\multirow{2}{*}{GLM-5} & \toolmemoryyes{} & \textbf{1.000} & \textbf{0.488} & \textbf{195.9} & \textbf{0.344$\times$} \\
 & \toolmemoryno{} & 0.889 & 0.434 & 500.9 & 1.000$\times$ \\
\midrule
\multirow{2}{*}{Gemini 3.1 Pro} & \toolmemoryyes{} & 0.889 & \textbf{0.469} & \textbf{309.9} & \textbf{0.137$\times$} \\
 & \toolmemoryno{} & 0.889 & 0.201 & 968.2 & 1.000$\times$ \\
\midrule
\multirow{2}{*}{Overall} & \toolmemoryyes{} & \textbf{0.963} & \textbf{0.534} & \textbf{242.5} & \textbf{0.327$\times$} \\
 & \toolmemoryno{} & 0.815 & 0.310 & 609.5 & 1.000$\times$ \\
\bottomrule
\end{tabular*}
\end{table}




On the same generated decks, the DeepPresenter-style quality metrics in
Table~\ref{tab:profile_memory_v6_ppteval_main} check whether persona gains
remain compatible with ordinary presentation quality. The results show that
\methodname improves personalization alignment while maintaining competitive
PPT generation quality.

%\FloatBarrier



\subsubsection{Localized Revision}


In this diagnostic matched-pair modify setting, tool-memory injection is
associated with stronger closed-loop completion and post-edit verification while
reducing time to the first correct edit and core tool time
(Table~\ref{tab:tool_memory_table1_redesign}). Since each matched pair changes
only tool-memory injection, these process gains support the combined
tool-memory pathway for reliable multi-turn modification, while the W-L-T-NA
counts expose pair-level heterogeneity; Appendix
Tables~\ref{tab:tool_memory_table1_pair_detail} and
\ref{tab:tool_memory_paired_robustness} provide pair-level details and the
paired robustness check.

The same locality-sensitive setting is illustrated qualitatively in
Figure~\ref{fig:localized_modify_example}, which contrasts a broader edit
footprint with a targeted patch on the requested element; Appendix
Figure~\ref{fig:appendix_tool_memory_localized_editing} shows a complementary
process trajectory.

\begin{figure}[!t]
  \centering
  \includegraphics[width=\linewidth]{localized_modify_example}
  \vspace{-4mm}
  \caption{Qualitative illustration of localized modify execution. Given the
  same local edit request, DeepPresenter can satisfy the target change while
  altering non-target regions of the slide. \methodname instead applies a
  targeted patch to the requested element, preserving already aligned slide
  content.}
  \label{fig:localized_modify_example}
\end{figure}

\FloatBarrier

\subsection{Analysis and Discussion}

\noindent\textbf{User profile memory yields broad, not marginal, alignment gains.}
Table~\ref{tab:profile_memory_v6_bestof_main} shows a consistent but not
uniform pattern across model families. With GLM-5 and Gemini 3.1 Pro,
\methodname leads both baselines on all four persona-alignment dimensions. With
GPT-5, the gains are strongest on Content and Specificity, while the two
baselines retain isolated advantages on Structure and Visual. Averaged across
model families, \methodname improves over DeepPresenter by 1.37 points on
Content, 0.53 on Structure, 1.66 on Visual, and 1.19 on Specificity, and
improves over SlideTailor by 2.73, 2.95, 2.79, and 3.08 points on the same
dimensions.

\noindent\textbf{The gains are planning-level persona gains.}
The strongest evidence is the joint movement of Structure and Specificity.
Structure excludes template retrieval accuracy, while Specificity uses
distractor personas; moreover, judges do not see the original prompt or parsed
intent. Thus the improvement cannot be read as merely better template matching
or generic visual polish. It indicates that routed long-term profiles help
decide page roles, ordering, layout fit, evidence emphasis, and persona-distinct
framing before round-0 generation.

\noindent\textbf{Persona gains remain compatible with general deck quality.}
Table~\ref{tab:profile_memory_v6_ppteval_main} shows that \textsc{Ours} obtains
the best Avg. on GPT-5 and remains competitive on GLM-5, while Gemini 3.1 Pro
shows the best Style and Diversity but lower Constraint. These results do not
claim uniform dominance on every quality metric; they show that the
persona-alignment improvements in Table~\ref{tab:profile_memory_v6_bestof_main}
are not simply a trade-off against ordinary presentation quality.

\begin{figure}[!t]
  \centering
  \includegraphics[width=\linewidth]{six_round_profile_consolidation_cases}
  \vspace{-2mm}
  \caption{Qualitative cross-job profile consolidation. Across six repeated
  jobs, local feedback cues become reusable profile preferences and later
  reappear as default slide-organization patterns, including evidence-boundary
  guardrails, owner/timeline closure tables, module input--output responsibility
  schemas, and implementation checklists. The figure is diagnostic case evidence
  rather than a separate quantitative metric.}
  \label{fig:six_round_profile_consolidation}
\end{figure}

\noindent\textbf{Localized revision should be judged by locality as well as success.}
Figure~\ref{fig:localized_modify_example} shows the qualitative difference
behind the modify setting: a system may satisfy the requested change while
still touching non-target regions, whereas \methodname keeps the patch closer
to the requested element. This is why the localized-revision metrics are
paired with process constraints such as closed-loop completion, verification,
and core-tool-time efficiency.

\noindent\textbf{Scope of profile and working-memory evidence.}
The round-0 persona-alignment table is the main quantitative evidence for user
profile memory. Figure~\ref{fig:six_round_profile_consolidation} provides
qualitative cross-job profile-consolidation evidence, showing how repeated local
cues become reusable organization preferences in later jobs. Appendix
Figure~\ref{fig:appendix_working_memory_carryover} provides complementary
within-session evidence for delayed carryover of active temporary memory, and
Appendix Figure~\ref{fig:appendix_template_generation_showcase} shows
template-guided generation examples.


\noindent\textbf{Tool memory improves overall reliability and search efficiency, but not by winning every pair.}
In the diagnostic matched-pair modify setting, tool-memory injection raises
overall Closed-Loop Completion from 0.815 to 0.963 (W-L-T-NA 3-1-5-0) and
Strict Verify from 0.310 to 0.534 (8-1-0-0). It also reduces Time to First
Correct Edit from 609.5s to 242.5s (6-2-0-1) and lowers the geometric mean
Core Tool Time Ratio to 0.327$\times$ (8-1-0-0), indicating less
non-inspection tool work. The pair-level counts are important: one Gemini
hard-modify pair loses on Closed-Loop Completion, two pairs lose on first-edit
latency, and one GPT-5 pair uses more core tool time. Together with the
pair-level detail in Appendix~\ref{app:tool_memory_pair_details}, these process
metrics are consistent with a more constrained localized editing pattern rather
than repeated full-deck regeneration or broad exploratory edits.

\FloatBarrier
